\documentclass[11pt, a4paper]{article}

\usepackage[a4paper, top=2.5cm, bottom=2.5cm, left=2cm, right=2cm]{geometry}
\usepackage{amsmath}
\usepackage{amssymb}
\usepackage[colorlinks=true, linkcolor=blue, urlcolor=blue, citecolor=blue]{hyperref}

\title{Theory of Everything - Volume 29 \\ \Large Ecosystem Dynamics \& Symbiosis}
\author{Omega Protocol}
\date{\today}

\begin{document}

\maketitle

\section*{Layman's Summary}
A forest is more than just a collection of trees; it's a single, massively interconnected super-organism. Volume 29 proves that competition and "survival of the fittest" eventually give way to cooperation and symbiosis, driven by the universe's relentless demand for thermodynamic efficiency.

\section{Symbiotic Integration (Axiom 32)}
While individuals compete, entire networks tend to merge. We establish that competing informational networks will inevitably combine into higher-order cooperative structures. Just as ancient bacteria merged to create the complex cells of plants and animals, ecosystems merge to minimize waste and internal friction. 

\section{Trophic Cascades (Theorem)}
If you remove wolves from a park, the rivers might change course. This is a trophic cascade. We prove mathematically that ecosystems are non-linear fractal networks. Changing or removing a top-level hub (like an apex predator) doesn't just affect its immediate prey; it sends mathematical shockwaves that reorganize the entire physical structure of the environment.

\section{Keystone Species (Theorem)}
In a stone arch, removing the central "keystone" causes the whole thing to collapse. We prove that ecosystems possess similar topological keystones. Certain biological nodes have such disproportionate importance to the flow of information and energy that their deletion causes a mathematical collapse of the local environmental manifold.

\end{document}
